
% =============================================================================
% GEODERMA MANUSCRIPT TEMPLATE
% Elsevier elsarticle class
% =============================================================================

\documentclass[preprint,12pt]{elsarticle}

% -----------------------------------------------------------------------------
% 1. PACKAGES
% -----------------------------------------------------------------------------

% Mathematics
\usepackage{amsmath}
\usepackage{amssymb}
\usepackage{mathtools}

% Figures
\usepackage{graphicx}
\usepackage{float}
\usepackage{subcaption}

% Tables
\usepackage{booktabs}
\usepackage{tabularx}
\usepackage{array}
\usepackage{multirow}
\usepackage{threeparttable}

% SI units
\usepackage{siunitx}
\sisetup{
    detect-all,
    separate-uncertainty = true
}

% Text formatting
\usepackage{lineno}
\usepackage{microtype}

% Hyperlinks and cross-referencing
\usepackage[colorlinks=true,
            linkcolor=blue,
            citecolor=blue,
            urlcolor=blue]{hyperref}

\usepackage[nameinlink,noabbrev]{cleveref}

% -----------------------------------------------------------------------------
% 2. JOURNAL AND BIBLIOGRAPHY SETTINGS
% -----------------------------------------------------------------------------

% Elsevier bibliography style
% Use numerical citations: [1], [2], [3]
\bibliographystyle{elsarticle-num}

% -----------------------------------------------------------------------------
% 3. JOURNAL METADATA
% -----------------------------------------------------------------------------

\journal{Geoderma}

% -----------------------------------------------------------------------------
% TITLE AND AUTHORS
% -----------------------------------------------------------------------------
\begin{document}
\begin{frontmatter}

\title{Feedstock-dependent effects of enhanced rock weathering on soil organic matter dynamics and persistence}

\author[aff1]{Teneille Nel\corref{cor1}}
\ead{first.author@email.com}

\author[aff1]{Lucilla Boito}
\ead{second.author@email.com}

\author[aff1]{Laura Steinwidder}

\author[aff2]{Carlos A. Sierra}

\author[aff1]{Sara Vicca}

\address[aff1]{
    Department of Soil Science,
    University Name,
    City,
    Country
}

\address[aff2]{
    Department of Environmental Science,
    University Name,
    City,
    Country
}

\cortext[cor1]{Corresponding author}

\begin{abstract}
Write your abstract here.
\end{abstract}

\begin{keyword}
Soil organic carbon
\sep Enhanced rock weathering
\sep Soil geochemistry
\sep Carbon persistence
\end{keyword}

\end{frontmatter}
% -----------------------------------------------------------------------------
% 7. MAIN TEXT
% -----------------------------------------------------------------------------

% Uncomment for line numbering during manuscript review
% \linenumbers

\section{Introduction}
\label{sec:introduction}

Significant increases in atmospheric CO$_2$ concentrations since the Industrial Revolution necessitate effective carbon dioxide removal (CDR) technologies. Enhanced rock weathering (ERW), which involves the application of crushed silicate rocks to soils, has the potential to contribute to long-term atmospheric CO$_2$ removal by accelerating natural silicate weathering \citep{} \citep{hartmann2008,manning2024}. During this process, atmospheric CO$_2$ is converted into dissolved bicarbonate, which can subsequently contribute to inorganic carbon storage through carbonate formation, as illustrated by the following reaction:

\begin{equation}
    \mathrm{CaSiO_3 + CO_2 + 2H_2O
    \rightarrow CaCO_3 + H_4SiO_4}
    \label{eq:carbonation}
\end{equation}

Agricultural soils are particularly suitable for ERW and soil carbon sequestration (SCS) due to the extensive land area available and the potential synergies with existing agricultural management practices. Although ERW research has traditionally focused on inorganic carbon (C) sequestration, recent studies have demonstrated that silicate rock amendments can also substantially influence soil organic matter (SOM) dynamics, potentially affecting organic C fluxes several orders of magnitude greater than those associated with soil inorganic carbon (SIC) fluxes. These effects may result in either increases or losses of SOM stocks \citep{vicca2022,steinwidder2025,sokol2024}. Understanding the effects of ERW on SOM dynamics is therefore essential for accurate net C accounting \citep{niron2024}, particularly because increases in soil organic carbon (SOC) stocks may generate additional carbon credits independently of those associated with ERW-driven SIC sequestration \citep{manning2024}.

The effects of ERW on SOM are governed by competing mechanisms that may either promote SOC stabilization or accelerate its mineralization. The release of divalent cations during silicate weathering can enhance organo-mineral associations through electrostatic cation bridging between organic matter and clay mineral surfaces. Furthermore, the formation of secondary minerals, including Al/Fe (oxy)(hydr)oxides, may provide additional sorption surfaces for SOM stabilization as mineral-associated organic matter (MAOM) \citep{vienne2025}. Conversely, increases in soil pH and the release of micronutrients may stimulate microbial activity and SOM decomposition. Enhanced production of organic acids by plant roots and microorganisms during nutrient acquisition from rock amendments may also destabilize existing mineral-organic associations, promoting the desorption of previously protected organic C and exposing it to microbial decomposition \citep{niron2024,sokol2024}. Additionally, increases in plant productivity and organic matter inputs may stimulate decomposition losses through enhanced microbial activity \citep{stoner2021}. Consequently, the net effect of ERW on SOC stocks depends not only on the magnitude of C inputs and losses but also on the mechanisms governing SOM formation and persistence.

\subsection{SOM persistence and carbon stability pools}

The persistence of SOC is fundamentally governed by its susceptibility to microbial decomposition and the environmental and physicochemical conditions that regulate its accessibility and stabilization. Six broad categories of controls on SOM persistence were recently summarized by \citep{rocci2024}: physical inaccessibility, organo-mineral and -metal stabilization, biodegradability of plant inputs, abiotic environmental factors (climate and soil chemistry), biochemical reactivity and diversity (recalcitrance), and microbial physiology and morphology. Given the substantial heterogeneity of SOM, separating it into distinct components that differ in their formation, function, and persistence is necessary to accurately predict the vulnerability of bulk SOC to environmental change \citep{lavallee2019}.

In particular, particulate organic matter (POM) and MAOM represent functionally distinct SOC pools with different stabilization mechanisms. POM persistence is primarily associated with the biochemical characteristics of organic matter and its physical protection within soil aggregates, whereas MAOM generally comprises relatively simple molecular compounds stabilized through interactions with mineral surfaces. These interactions include van der Waals forces, hydrogen bonding, covalent bonding, complexation, cation bridging, and pore sorption. Distinguishing between these SOC fractions is therefore essential for identifying how ERW influences the formation and persistence of different organic C pools and for evaluating the mechanisms responsible for changes in bulk SOC stability.

\subsection{Assessing SOC persistence using radiocarbon}

Quantifying both net C inputs and the transit time of C in soils is essential for evaluating the effectiveness of soil C sequestration strategies in mitigating climate change \citep{sierra2021,xia2013}. Although the effects of ERW on SOC accrual are increasingly being investigated \citep{boito2025,steinwidder2025}, the persistence of ERW-associated SOC remains poorly constrained. Assigning a precise stabilization time to SOC is particularly challenging and requires further research \citep{manning2024}. This uncertainty is especially relevant for topsoil SOM, which is in direct contact with the atmosphere and is consequently particularly vulnerable to decomposition \citep{mills2014}.

While fast-cycling SOC pools can be investigated through relatively short-term experiments, the majority of soil C is stored in slower-cycling pools that are less amenable to experimental investigation \citep{mills2014}. Radiocarbon ($^{14}$C) measurements, combined with observed SOC stocks, provide a means of constraining models of SOM dynamics over longer timescales, assuming steady-state conditions \citep{sierra2014}. Naturally occurring $^{14}$C decay can be used to estimate centennial-scale turnover, whereas the anthropogenic bomb-$^{14}$C spike associated with atmospheric nuclear testing during the mid-twentieth century provides information on SOC dynamics over decadal timescales \citep{mills2014}. Integrating radiocarbon measurements with SOC fractionation and modelling approaches can therefore provide insights into the persistence of distinct SOC pools and their responses to ERW-induced changes in soil properties.

\subsection{Soil pH as a driver of SOC stabilization and destabilization}

Soil pH is a major determinant of the mechanisms governing SOC stabilization and decomposition. As pH increases, the dominant mechanisms of SOC stabilization may shift from organo-metal complexation, primarily involving carboxyl coordination and hydrogen bonding, towards associations with short-range-order (SRO) mineral phases and, subsequently, Ca complexation and cation bridging with phyllosilicates \citep{zhou2024,rasmussen2018,jiang2026}. Increasing pH promotes the deprotonation of carboxyl groups, potentially enhancing electrostatic repulsion and destabilizing mineral-organic associations, thereby decreasing MAOM accumulation \citep{jiang2026}. In addition, soil pH regulates microbial diversity and activity and is consequently an important driver of SOM decomposition \citep{wangkuzyakov2024}.

Evidence from ERW experiments suggests that pH-induced changes in mineral-organic interactions may result in SOC destabilization. For example, \citep{steinwidder2025} reported evidence of SOC release from MAOM following an increase in soil pH associated with silicate rock powder amendments. However, the consequences of increased rhizosphere respiration for SOC stocks are not necessarily negative, as enhanced respiration is often accompanied by increased photosynthesis and C inputs, potentially offsetting decomposition losses \citep{kuzyakovcheng2001}. Moreover, roots and rhizodeposits represent important precursors of MAOM and contribute to aggregate formation \citep{villarino2021,amelung2023}. Organic matter may also facilitate mineral nucleation and transformation, further influencing the formation and persistence of organo-mineral associations \citep{kleber2021}. Thus, the net effect of ERW-induced pH changes on SOC persistence likely depends on the balance between destabilization of existing associations, enhanced microbial decomposition, and the formation of new mineral-organic associations.

\subsection{Secondary mineral formation and SOC stabilization}

ERW can accelerate mineral dissolution beyond natural weathering rates, increasing soil solution solute concentrations and potentially reducing primary mineral dissolution rates while stimulating the formation of secondary minerals \citep{bluem2016,steinwidder2025}. The subglacial grinding of bedrock has been shown to produce amorphous Si, Al, and Fe phases that may act as precursors to more crystalline phases, including phyllosilicates \citep{rampe2022,behrens2015}. The formation of secondary clay minerals during ERW may subsequently influence weathering rates and reduce the efficacy of inorganic CDR \citep{manning2024}. However, the formation of reactive SRO metal oxides, such as allophane and ferrihydrite, and aluminosilicates, such as proto-imogolite, may provide additional mineral surfaces for SOM sorption and stabilization \citep{sokol2024,steinwidder2025}. These processes could therefore enhance MAOM formation and promote long-term SOM stabilization in young soils.

The interactions between SOM and secondary minerals may also influence mineral formation itself. Accumulation of SOM on mineral surfaces can facilitate the nucleation of new mineral particles, inhibit the growth of existing particles, and promote the formation and persistence of smaller SRO mineral phases. In this context, \citep{steinwidder2025} observed that base cations released during enhanced silicate weathering were not recovered in sequentially extracted pools, suggesting their potential incorporation into more stable mineral pools, such as clays or other secondary phases. However, the extent to which ERW-induced secondary mineral formation contributes to SOC stabilization, and how these processes influence SOC persistence, remains uncertain.

The mineralogical characteristics of secondary phases are also likely to determine their effectiveness in stabilizing SOC. Chemisorption is considered a dominant control on C turnover in more weathered soils containing greater proportions of crystalline aluminosilicates and Fe and Al (oxy)hydroxides, such as goethite and hematite \citep{khomo2017}. These minerals may stabilize SOC more effectively than non-crystalline phases due to stronger ligand-exchange bonds \citep{sun2026}. In a subhumid region, bulk soil C turnover time was strongly correlated with smectite content, reflecting the high charge density and surface area of 2:1 clays and their consequently high affinity for C. In contrast, the overall C inventory was more closely related to crystalline Fe and Al oxyhydroxides \citep{khomo2017}. Similarly, \citep{sun2026} found that amorphous allophane exhibited greater C storage capacity due to its higher specific surface area (SSA), whereas crystalline kaolinite demonstrated stronger inner-sphere binding with SOM through direct coordination between organic matter hydroxyl groups and mineral surface siloxane groups \citep{sun2026}. These findings highlight the importance of distinguishing between mineral phases and their respective stabilization mechanisms when assessing the potential effects of ERW on SOC persistence.

\subsection{Cation-mediated stabilization and soil aggregation}

Mineral surface charge is another important determinant of SOC stabilization and soil aggregation. Minerals exhibit both permanent structural charges and variable charges that depend on pH and adsorbed ions. These charge distributions regulate processes including coagulation, dispersion, colloidal stability, solubility, and adsorptive bonding \citep{kleber2021}. Through silicate weathering, ERW can release divalent cations, such as Ca$^{2+}$, which may enhance SOM stabilization through cation bridging, co-precipitation with carbonate minerals, and the production of microbial biofilms \citep{sokol2024,rowley2018,buss2024}. Similarly, the addition of non-structural Fe ions to allophanes has been observed to increase SOC persistence through cation bridging effects \citep{sun2026}.

Cation composition and electrolyte concentrations also influence clay flocculation and soil aggregation. The threshold electrolyte concentrations and pH required for clay flocculation vary according to soil mineralogy, with the ratio of divalent to monovalent cations and the net charge of clay minerals playing important roles. Clay dispersion in response to adsorbed cations is dependent on the net charge available for clay-water interactions \citep{marchuk2013}. Consequently, ERW-induced changes in cation concentrations, soil pH, and mineral surface properties may alter aggregate formation and stability, potentially influencing the physical protection and persistence of SOC. However, the relative contributions of cation-mediated organo-mineral interactions, secondary mineral formation, and aggregate stabilization to ERW-induced changes in SOC persistence remain insufficiently understood.

\subsection{Relevance of the Westmalle field trial}

Despite growing interest in ERW and increasing efforts to evaluate the CDR potential and associated co-benefits of different silicate feedstocks, field-based comparisons of materials with contrasting mineralogical compositions remain scarce \citep{rijnders2025,steinwidder2025,tepas2023}. Furthermore, relatively few studies have investigated silicate amendments alongside conventional agricultural practices, such as liming \citep{araujo2024,buckinghamhenderson2024,dietzen2018}. Evidence for ERW-induced agronomic benefits has predominantly been obtained from acidic, highly weathered, or nutrient-limited soils, whereas relatively fertile agricultural soils, which may be targeted for ERW deployment primarily for CDR rather than agronomic improvement, have received comparatively less attention and generally exhibit weaker responses \citep{swoboda2022,vanstraaten2002}. These knowledge gaps are also relevant to SOC persistence, as the effects of silicate amendments on organo-mineral interactions, secondary mineral formation, and aggregate stabilization may differ depending on initial soil conditions and feedstock mineralogy.

The Westmalle agricultural field trial in Belgium provides an opportunity to investigate these mechanisms under field conditions in a relatively fertile agricultural soil. The experimental site is characterized by a loamy sand Gleyic Podzol with near-neutral pH, relatively high SOC content, and moderate CEC. The trial compares two basalts (Hühnerberg and Bolsdorfer), a tephrite (Eifelgold), agricultural lime, and an unamended control, allowing the effects of silicate feedstocks with contrasting mineralogical compositions to be evaluated alongside conventional liming. While previous investigations at this site have focused on soil water dynamics, chemical fertility, and crop performance, the effects of these amendments on SOC persistence and the mechanisms governing the stabilization of distinct SOM fractions remain insufficiently constrained. In particular, it is unclear whether differences in feedstock weathering potential and mineralogical composition translate into contrasting effects on SOC persistence through changes in soil geochemistry, secondary mineral formation, cation-mediated organo-mineral associations, and aggregate stability.

\subsection{Knowledge gaps, hypotheses, and objectives}

Despite growing evidence that ERW influences SOC dynamics, the mechanisms governing the stabilization and destabilization of SOM under ERW remain poorly constrained. In particular, the extent to which ERW-induced changes in soil pH, secondary mineral formation, cation release, and soil aggregation affect the persistence of distinct SOC fractions is uncertain. Furthermore, the relative importance of these abiotic drivers in controlling SOC persistence across different stability pools remains insufficiently quantified. Addressing these knowledge gaps is necessary to determine whether ERW-induced changes in soil properties promote long-term SOC stabilization or increase the vulnerability of existing SOC to mineralization.

This study investigates ERW-induced mechanisms of SOM stabilization and destabilization, including organo-mineral associations and physical occlusion, and provides a suite of SOC persistence metrics for an ongoing Belgian ERW field trial. We hypothesize that ERW-induced secondary mineral formation and cation release promote soil aggregation and SOC binding, thereby enhancing SOC stabilization. Conversely, we hypothesize that increases in soil pH associated with silicate weathering may promote SOC mineralization and SOM desorption, decreasing the stability of both labile fPOM and MAOM.

The objectives of this study are to (i) quantify the persistence of distinct SOM stability fractions and (ii) identify the abiotic drivers governing SOC stability. By integrating SOC persistence metrics with geochemical and mineralogical indicators, this study aims to improve understanding of the mechanisms controlling SOC stabilization and destabilization under ERW and to provide insights into the potential implications for long-term soil C sequestration.

  
\section{Materials and methods}
\label{sec:methods}

\subsection{Study site}
\label{subsec:study_site}

Describe the study site.

\subsection{Experimental design}
\label{subsec:experimental_design}

Describe the experimental design and treatments.

\subsection{Soil analyses}
\label{subsec:soil_analyses}

Describe the analytical methods.
  The FeDCB/Fe total ratio relates total pedogenetically transformed Fe phases extracted with DCB, to total Fe, a value which increases with soil development as more crystalline oxide forms. (ref: maier 2025)

\subsection{Statistical analyses}
\label{subsec:statistics}

Describe the statistical analyses and models.

% Example mathematical equation
\begin{equation}
    SOC_{\mathrm{stock}} =
    C_{\mathrm{total}} \times
    \rho_{\mathrm{b}} \times
    z \times
    A
    \label{eq:soc_stock}
\end{equation}

where $C_{\mathrm{total}}$ is the soil carbon concentration,
$\rho_{\mathrm{b}}$ is bulk density, $z$ is soil depth,
and $A$ is the area conversion factor.

\section{Results}
\label{sec:results}

\subsection{Soil geochemical properties}
\label{subsec:geochemistry}

Present results here.

\subsection{Soil carbon fractions}
\label{subsec:carbon_fractions}

Present results here.
Bulk SOC changes are difficult to detect especially at lower sampling density due to soil heterogeneity (poeplau 2022, steinwidder 2025). 

\subsection{Relationships between soil properties}
\label{subsec:relationships}

Present results here.

\section{Discussion}
\label{sec:discussion}

\subsection{Effects of treatment on soil properties}
\label{subsec:treatment_effects}

The effects of ERW on soil properties provide important context for understanding the mechanisms underlying SOC stabilization and destabilization. The fate of base cations released during silicate weathering is particularly relevant, as their interactions with soil exchange complexes and secondary minerals may influence both soil chemistry and SOC persistence. \citep{steinwidder2025} found that the majority of base cations released during ERW were either sorbed to the exchangeable complex or incorporated into secondary minerals, such as (hydr)oxides and aluminosilicate clays, whereas only negligible amounts were recovered in pedogenic carbonates. This suggests that the effects of ERW on soil C dynamics may extend beyond inorganic C sequestration through carbonate formation, with potentially important implications for organo-mineral associations and SOC stabilization.

Despite these changes in soil chemistry, SOC stocks remained unaffected by silicate amendment in a 15-month mesocosm experiment \citep{steinwidder2025}. The authors attributed this response, at least partly, to increased plant C inputs that offset initial SOC losses through mineralization. This highlights the importance of considering the balance between C inputs and losses when interpreting changes in SOC stocks following ERW, as unchanged SOC stocks do not necessarily imply an absence of underlying changes in SOM dynamics.

The relationships between soil fertility and SOC may also provide insights into the mechanisms governing SOC accumulation. \citep{maier2025} reported positive correlations between SOC content and soil fertility indicators, including cation exchange capacity (CEC) and exchangeable bases, whereas exchangeable K and soil pH were negatively correlated with SOC. Soil metal oxides were also positively correlated with SOC content, while poorly crystalline oxides exhibited negative correlations \citep{maier2025}. These contrasting relationships suggest that the effects of soil chemical properties on SOC are not uniform, but may depend on the specific mineral phases and cations involved, as well as their influence on organo-mineral interactions and microbial activity.

\subsection{Mechanisms of soil carbon stabilization}
\label{subsec:stabilization}

\subsubsection{MAOM formation and mineral saturation}

The capacity of mineral surfaces to stabilize organic C is constrained by the availability of sorption sites. A global dataset analyzed by \citet{zhou2024} supported the concept of MAOM saturation, whereby the capacity for further direct organic matter association with mineral surfaces becomes limited as available binding sites are occupied. The degree of mineral saturation may therefore influence the efficiency with which newly added organic matter is transferred into MAOM. A greater C deficit, reflecting a larger difference between the existing MAOM content and its potential stabilization capacity, has been associated with increased MAOM transfer efficiency from POM to MAOM and enhanced MAOM stabilization \citep{stewart2008}.

However, the relationship between mineral saturation and SOC stabilization is not necessarily linear. \citep{sun2026} proposed that organic matter loading beyond mineral saturation increases the proportion of dynamic, outer-sphere-associated OM within MAOM, potentially making this fraction more vulnerable to mineralization. Thus, the quantity of organic matter associated with mineral surfaces does not necessarily reflect its long-term persistence, as the nature and strength of mineral-organic interactions may change with increasing C loading.

The importance of these mechanisms under ERW is illustrated by \citep{sokol2024}, who observed reduced MAOM stocks and POM accrual in the topsoil of irrigated cropland fields following two years of crushed rock additions. These changes were partly attributed to reduced microbial biomass and decreased transformation of POM into MAOM. This finding suggests that ERW-induced changes in soil properties may influence not only the capacity of mineral surfaces to stabilize organic C, but also the microbial processes responsible for MAOM formation. Consequently, the response of SOC fractions to ERW may depend on both the availability of mineral stabilization sites and the efficiency of microbial-mediated organic matter transformation.

\subsubsection{Secondary mineral formation and organo-mineral stabilization}

The formation and transformation of secondary minerals may provide an important mechanism for SOC stabilization following ERW. In particular, the mineralogical composition and degree of crystallinity of newly formed phases may determine their capacity to associate with organic matter. Al$_{\mathrm{ox}}$:Si$_{\mathrm{ox}}$ molar ratios of 2--2.3 are consistent with the presence of SRO aluminosilicates, such as allophane and imogolite \citep{watanabe2023}. Such mineral phases may provide reactive surfaces for organic matter sorption and contribute to MAOM formation.

Evidence for the role of secondary mineral formation in SOC stabilization has been reported across different soil systems. \citep{steinwidder2025} found that SOC was stabilized within the amorphous mineral pool following steel slag application. Similarly, concentrations of pedogenic Fe, Al, and Mn oxides were positively correlated with SOC stocks and radiocarbon content in alpine soils \citep{maier2025}. These findings indicate that metal-bearing mineral phases may contribute to SOC accumulation and persistence, although their effectiveness likely depends on their mineralogical characteristics and the nature of their interactions with organic matter.

An ERW study by \citep{matsuoka2025} provided further evidence for mineral-mediated SOC stabilization by demonstrating the co-localization of C, Fe, and Al within stable micro-aggregates. This was interpreted as evidence of a basalt-induced organo-metal glue effect, which was more pronounced in finer basalts characterized by faster weathering rates and higher concentrations of reactive metals. The authors suggested that the associated MAOM could form through co-precipitation, whereby SOM is adsorbed onto Fe-rich secondary minerals, or through complexes involving released metal(loid) cations \citep{matsuoka2025}. These mechanisms highlight the potential for ERW to promote SOC stabilization through the formation of new mineral-organic associations rather than solely through the sorption of organic matter onto pre-existing mineral surfaces.

Nevertheless, the persistence of these associations may depend on the stability of the mineral phases themselves. In soils experiencing strong wet-dry cycling, metastable minerals may be comparatively transient and rapidly transform into kaolinite. Consequently, the strength of C association with mineral phases alone may not adequately predict SOC persistence; the timescale over which these mineral phases remain stable within a given pedogenetic environment must also be considered \citep{khomo2017}. This distinction is particularly relevant to ERW, where the formation of reactive secondary minerals may initially enhance SOC stabilization, but subsequent mineral transformation could alter their C storage capacity over longer timescales.

\subsubsection{Soil aggregation and cation-mediated stabilization}

In addition to direct mineral-organic associations, ERW may influence SOC persistence through changes in soil aggregation and the physical occlusion of organic matter. \citep{steinwidder2025} found that silicate amendment stimulated macro-aggregate formation and the occlusion of SOC within these aggregates, suggesting that ERW may enhance physical protection of organic C. The formation and stabilization of aggregates are closely linked to the nature of organic matter inputs and their interactions with mineral particles. \citep{verchot2011} found evidence that macro-aggregates form around relatively fresh plant-derived organic matter particles characterized by carboxylic and aromatic functional groups produced through lignin oxidation. In contrast, micro-aggregates are formed through interactions between minerals and organic ligands, including microbially derived polysaccharides and aliphatic C associated with carbohydrate decomposition \citep{verchot2011}. These differences suggest that aggregate formation may contribute to SOC stabilization through distinct mechanisms operating at different spatial and temporal scales.

Calcium-mediated aggregation is particularly relevant in the context of ERW, given the release of Ca$^{2+}$ during silicate weathering. Higher Ca contents can increase soil aggregation and the amount of SOC occluded within aggregates. In humid environments, this is generally attributed to the flocculation of soil particles by exchangeable Ca$^{2+}$, whereas in arid regions, aggregation may occur through cementation associated with the evaporative precipitation of CaCO$_3$ \citep{muneeroades1989,fernandezugalde2014}. Furthermore, \citep{rowley2021} proposed that Ca increased MAOM-C to a greater extent than POM-C in Swiss rangeland soils, suggesting that Ca-mediated stabilization may extend beyond physical aggregation to include direct mineral-organic interactions.

The effects of Ca on SOM stabilization may also depend on soil pH and mineral surface charge. At higher pH, the negative variable surface charge of poorly crystalline metal oxide minerals increases, potentially promoting their interactions with cations. Accordingly, Ca$^{2+}$ addition may enhance SOM sorption through the formation of mineral-Ca-OM ternary complexes \citep{rowley2021,jiang2026}. Similar mechanisms have been proposed for Fe, with the addition of Fe in an experiment by \citep{sun2026} suggested to enhance SOM stabilization through mineral-Fe-SOM ternary complexes. These findings indicate that the effects of ERW-induced cation release on SOC persistence may involve both aggregate-mediated physical protection and enhanced organo-mineral binding.

\subsubsection{The role of Mn in SOC stabilization and decomposition}

Transition metal ions, including Mn, Fe, Ni, Cu, and Zn, can influence the electronic states of organic functional groups in solution or at mineral-water interfaces, potentially facilitating direct covalent bonding through electron transfer \citep{kleber2021}. In particular, Mn-(oxy)hydroxides may contribute to SOM stabilization through the formation of inner-sphere complexes, including ligand-exchange interactions. However, Mn may also influence SOC persistence indirectly through its effects on microbial nutrient availability and decomposition processes.

Metal oxide minerals can adsorb phosphates \citep{franks2021}, potentially reducing the availability of P for microbial growth. Such mineral sorption-induced P limitation has been shown to inhibit microbial decomposition activity, thereby increasing C persistence \citep{bramble2025}. This suggests that metal oxide formation may contribute to SOC stabilization not only through direct sorption of organic matter but also through changes in microbial nutrient availability.

In contrast, Mn can also promote SOM decomposition through its redox activity. Mn$^{2+}$ is the bioavailable form of Mn that can be oxidized by fungal enzymes, whereas Mn$^{4+}$ is a strong oxidant that has been shown to oxidize SOM \citep{paulus2024}. Several studies have reported negative correlations between Mn availability and SOC storage, attributed to increased SOC decomposition \citep{berg2000,keiluweit2015}. Consequently, the role of Mn in SOC dynamics is potentially dual, depending on its mineralogical form, oxidation state, and interactions with microbial processes. While Mn-bearing mineral phases may contribute to organo-mineral stabilization, bioavailable Mn and its redox transformations may also enhance SOM oxidation and decomposition. The net effect of Mn on SOC persistence may therefore depend on the balance between these competing stabilization and destabilization mechanisms.

\subsection{Implications for enhanced rock weathering}
\label{subsec:erw_implications}

% To be developed following the integration of experimental results.
% Discuss the implications of the observed treatment effects for:
% (i) SOC persistence and the stabilization of distinct SOM fractions;
% (ii) the relative importance of secondary mineral formation, cation bridging,
% and aggregate-mediated physical protection;
% (iii) potential trade-offs between SOC stabilization and mineralization;
% and (iv) the implications for ERW carbon accounting and long-term CDR.

\section{Conclusions}
\label{sec:conclusions}

Summarize the main findings and their implications.

% -----------------------------------------------------------------------------
% 8. ACKNOWLEDGEMENTS
% -----------------------------------------------------------------------------

\section*{Acknowledgements}

Acknowledge funding, technical assistance, and collaborators.

% -----------------------------------------------------------------------------
% 9. DECLARATIONS
% -----------------------------------------------------------------------------

\section*{Declaration of competing interest}

The authors declare that they have no known competing
financial interests or personal relationships that could
have appeared to influence the work reported in this paper.

\section*{Data availability}

The data supporting the findings of this study are available
from the corresponding author upon reasonable request.

% -----------------------------------------------------------------------------
% 10. REFERENCES
% -----------------------------------------------------------------------------

\bibliography{references}

% -----------------------------------------------------------------------------
% 11. SUPPLEMENTARY MATERIAL
% -----------------------------------------------------------------------------

% If supplementary material is submitted as a separate file,
% reference it in the manuscript as appropriate.

\end{document}
